The ProcessSet class is a set container for multiple
Process objects. It shares many of the same operations as the Process class,
but when an operation is performed on a ProcessSet it is done on every Process
in the ProcessSet. On some systems, a ProcessSet can
achieve better performance when repeating an operation across many target
processes. It satisfies the C++
\href{https://en.cppreference.com/cpp/named_req/ForwardIterator}{LegacyForwardIterator}
concept. It also offers the standard \code{begin/end} members for iteration.

\begin{apient}
  ProcessSet::ptr
  ProcessSet::const_ptr
\end{apient}

\apidesc{
  The ptr and const\_ptr types are smart pointers to a ProcessSet object.
}

\begin{apient}
  struct CreateInfo
  struct AttachInfo
\end{apient}

\apidesc{
  The CreateInfo and AttachInfo types are used by the
  ProcessSet::createProcessSet and
  ProcessSet::attachProcessSet functions when creating groups of processes.
}

\begin{apient}
  struct write_t
  struct read_t
\end{apient}

\apidesc{
  The write\_t and read\_t types are used by ProcessSet::readMemory and
  ProcessSet::writeMemory.
}

\begin{apient}
  static ProcessSet::ptr newProcessSet()
\end{apient}

\apidesc{
  This function creates a new ProcessSet that is empty.
}

\begin{apient}
  static ProcessSet::ptr newProcessSet(Process::const_ptr proc)
\end{apient}

\apidesc{
  This function creates a new ProcessSet containing proc.
}

\begin{apient}
  static ProcessSet::ptr newProcessSet(ProcessSet::const_ptr pset)
\end{apient}

\apidesc{
  This function creates a new ProcessSet that is a copy of pset.
}

\begin{apient}
  static ProcessSet::ptr newProcessSet(const std::set<Process::const_ptr> &procs)
\end{apient}

\apidesc{
  This function creates a new ProcessSet containing every element from procs.
}

\begin{apient}
  static ProcessSet newProcessSet(AddressSet::const_iterator ab, AddressSet::const_iterator ae)
\end{apient}

\apidesc{
  This function creates a new ProcessSet containing the processes that are found
  within \code{[ab, ae)} of an AddressSet.
}

\begin{apient}
  static ProcessSet::ptr createProcessSet(std::vector<CreateInfo> &cinfo)
\end{apient}

\apidesc{
  This function creates a new ProcessSet by launching new processes. Each
  element in cinfo specifies an executable, arguments, environment and file
  descriptor mappings (with similar semantics to
  Process::createProcess), which are used
  to launch a new process.
  Every successfully created Process will be added to a new ProcessSet that is
  returned by this function.
  In addition, the cinfo vector will be updated so that each
  entry\' s proc field points to the Process created by that
  entry, and the error\_ret entry will contain an error code for any process
  launch that failed.
}

\begin{apient}
  static ProcessSet::ptr attachProcessSet(std::vector<AttachInfo> &ainfo)
\end{apient}

\apidesc{
  This function creates a new ProcessSet by attaching to existing processes.
  Each element in ainfo specifies a PID and executable (with similar semantics
  to Process::attachProcess), which are
  used to attach to the processes.
  Every successfully attached Process will be added to a new ProcessSet that is
  returned by this function.
  In addition, the ainfo vector will be updated so that each
  entry\' s proc field points to the Process attached by that
  entry, and the error\_ret entry will contain an error code any process attach
  that failed.
}

\begin{apient}
  ProcessSet::ptr set_union(ProcessSet::ptrpset) const
\end{apient}

\apidesc{
  This function returns a new ProcessSet whose elements are a set union of this
  ProcessSet and pset.
}

\begin{apient}
  ProcessSet::ptr set_intersection(ProcessSet::ptr pset) const
\end{apient}

\apidesc{
  This function returns a new ProcessSet whose elements are a set
  intersection of
  this ProcessSet and pset.
}

\begin{apient}
  ProcessSet::ptr set_difference(ProcessSet::ptr pset) const
\end{apient}

\apidesc{
  This function returns a new ProcessSet whose elements are a set difference of
  this ProcessSet minus pset.
}

\begin{apient}
  iterator find(Process::const_ptr proc)
  const_iterator find(Process::const_ptr proc) const
\end{apient}

\apidesc{
  These functions search a ProcessSet for the Process pointed to by proc and
  returns an iterator that points to that element. It returns ProcessSet::end()
  if no element is found.
}

\begin{apient}
  iterator find(Dyninst::PID pid)
  const_iterator find(Dyninst::PID pid) const
\end{apient}

\apidesc{
  These functions search a ProcessSet for the Process pointed to by proc and
  returns an iterator that points to that element. It returns ProcessSet::end()
  if no element is found.
}

\begin{apient}
  bool empty() const
\end{apient}

\apidesc{
  This function returns true if the ProcessSet has zero elements,
  false otherwise.
}

\begin{apient}
  size_t size() const
\end{apient}

\apidesc{
  This function returns the number of elements in the ProcessSet.
}

\begin{apient}
  std::pair<iterator, bool> insert(Process::const_ptr proc)
\end{apient}

\apidesc{
  This function inserts proc into the ProcessSet. If proc already exists in the
  ProcessSet, then no change will occur. This function returns an iterator
  pointing to either the new or existing element and a boolean that
  is true if an
  insert happened and false otherwise.
}

\begin{apient}
  void erase(iterator pos)
\end{apient}

\apidesc{
  This function removes the element pointed to by pos from the ProcessSet.
}

\begin{apient}
  size_t erase(Process::const_ptr proc)
\end{apient}

\apidesc{
  This function searches the ProcessSet for proc, then erases that element from
  the ProcessSet. It returns 1 if it erased an element and 0 otherwise.
}

\begin{apient}
  void clear()
\end{apient}

\apidesc{
  This function erases all elements in the ProcessSet.
}

\begin{apient}
  ProcessSet::ptr getErrorSubset() const
\end{apient}

\apidesc{
  This function returns a new ProcessSet containing every Process from this
  ProcessSet that has a non-zero error code. Error codes are reset upon every
  ProcessSet API call, so this function shows which Processes had an
  error on the
  last ProcessSet operation.
}

\begin{apient}
  void getErrorSubsets(std::map<ProcControlAPI::err_t, ProcessSet::ptr> &err_sets) const
\end{apient}

\apidesc{
  This function returns a set of new ProcessSets containing every Process from
  this ProcessSet that has non-zero error codes, and grouped by error code. For
  each error code generated by the last ProcessSet API operation an element will
  be added to err\_sets, and every Process that has the same error code will be
  added to the new ProcessSet associated with that error code.
}

\begin{apient}
  bool anyTerminated() const
  bool allTerminated() const
\end{apient}

\apidesc{
  These functions respectively return true if any or all processes in this
  ProcessSet are terminated, and false otherwise.
}

\begin{apient}
  bool anyExited() const
  bool allExited() const
\end{apient}

\apidesc{
  These functions respectively return true if any or all processes in this
  ProcessSet have exited normally, and false otherwise.
}

\begin{apient}
  bool anyCrashed() const
  bool allCrashed() const
\end{apient}

\apidesc{
  These functions respectively return true if any or all processes in this
  ProcessSet have crashed normally, and false otherwise.
}

\begin{apient}
  bool anyDetached()
  bool allDetached()
\end{apient}

\apidesc{
  These functions respectively return true if any or all processes in this
  ProcessSet have been detached, and false otherwise.
}

\begin{apient}
  bool anyThreadStopped()
  bool allThreadStopped()
\end{apient}

\apidesc{
  These functions respectively return true if any or all threads in this
  ProcessSet are stopped, and false otherwise.
}

\begin{apient}
  bool anyThreadRunning()
  bool allThreadRunning()
\end{apient}

\apidesc{
  These functions respectively return true if any or all threads in this
  ProcessSet are running, and false otherwise.
}

\begin{apient}
  ProcessSet::ptr getTerminatedSubset() const
\end{apient}

\apidesc{
  This function returns a new ProcessSet, which is a subset of this ProcessSet,
  and contains every Process that is terminated.
}

\begin{apient}
  ProcessSet::ptr getExitedSubset() const
\end{apient}

\apidesc{
  This function returns a new ProcessSet, which is a subset of this ProcessSet,
  and contains every Process that has exited normally.
}

\begin{apient}
  ProcessSet::ptr getCrashedSubset() const
\end{apient}

\apidesc{
  This function returns a new ProcessSet, which is a subset of this ProcessSet,
  and contains every Process that has crashed.
}

\begin{apient}
  ProcessSet::ptr getDetachedSubset() const
\end{apient}

\apidesc{
  This function returns a new ProcessSet, which is a subset of this ProcessSet,
  and contains every Process that is detached.
}

\begin{apient}
  ProcessSet::ptr getAllThreadRunningSubset() const
  ProcessSet::ptr getAnyThreadRunningSubset() const
\end{apient}

\apidesc{
  This function returns a new ProcessSet, which is a subset of this ProcessSet,
  and contains every Process that respectively has any or all threads running.
}

\begin{apient}
  ProcessSet::ptr getAllThreadStoppedSubset() const
  ProcessSet::ptr getAnyThreadStoppedSubset() const
\end{apient}

\apidesc{
  This function returns a new ProcessSet, which is a subset of this ProcessSet,
  and contains every Process that respectively has any or all threads stopped.
}

\begin{apient}
  bool continueProcs()
\end{apient}

\apidesc{
  This function continues every thread in every process of this ProcessSet,
  similar to Process::continueProc. It
  returns true if every process was successfully continued and false otherwise.
}

\begin{apient}
  bool stopProcs()
\end{apient}

\apidesc{
  This function stops every thread in every process of this
  ProcessSet, similar to
  Process::stopProc. It returns true if every process was
  successfully stopped and false otherwise.
}

\begin{apient}
  bool detach(bool leaveStopped = true)
\end{apient}

\apidesc{
  This function detaches from every process in this ProcessSet, similar to
  Process::detach. It returns true if every
  process detach was successful and false otherwise.
  If the leaveStopped parameter is set to true, and the processes in this
  ProcessSet are stopped, then the processes will be left in a stopped state
  after the detach.
}

\begin{apient}
  bool terminate()
\end{apient}

\apidesc{
  This function terminates every process in this ProcessSet, similar to
  Process::terminate. It returns true if
  every process was successfully terminated and false otherwise.
}

\begin{apient}
  bool temporaryDetach()
\end{apient}

\apidesc{
  This function does a temporary detach from every process in this ProcessSet,
  similar to Process::temporaryDetach.
  It returns true if every process was successfully detached and false
  otherwise.
}

\begin{apient}
  bool reAttach()
\end{apient}

\apidesc{
  This function reattaches to every process in this ProcessSet, similar to
  Process::reAttach. It returns true if every
  process was successfully reAttached and false otherwise.
}

\begin{apient}
  AddressSet::ptr mallocMemory(size_t sz) const
\end{apient}

\apidesc{
  This function allocates a block of memory of size sz in each process in this
  ProcessSet. The addresses of the allocations are returned in a new AddressSet
  object.
}

\begin{apient}
  bool mallocMemory(size_t size, AddressSet::ptr location)
\end{apient}

\apidesc{
  This function allocates a block of memory of size sz in each process in this
  ProcessSet. The memory will be allocated in each process based on the
  Process/Address pairs in location.
  This function\' s behavior is undefined if location contains
  processes not included in this ProcessSet.
  This function returns true if every allocation happened without
  error and false
  otherwise.
}

\begin{apient}
  bool freeMemory(AddressSet::ptr addrs) const
\end{apient}

\apidesc{
  This function frees memory allocated by
  Process::mallocMemory or
  ProcessSet::mallocMemory. The
  AddressSet addrs should contain a list of Process/Address pairs that point to
  the memory that should be freed.
  This function\' s behavior is undefined if addrs contains
  processes not included in this ProcessSet.
  This function returns true if every free happened without error and false
  otherwise.
}

\begin{apient}
  bool readMemory(AddressSet::ptr addrs,
                  std::multimap<Process::ptr, void *> &result,
                  size_t size) const
\end{apient}

\apidesc{
  This function reads memory from processes in this ProcessSet. addrs should
  contain the addresses to read memory from. size should be the amount of
  memory read from each process. The results of the memory reads will be
  returned by filling in the result multimap. Each process that is read from
  will have an entry in result along with a malloc allocated buffer containing
  the results of the read.
  It is the ProcControlAPI user\' s responsibility to free the
  memory buffers returned by this function.
  This function\' s behavior is undefined if addrs contains
  processes not included in this ProcessSet.
  This function returns true if every read happened without error, and false
  otherwise.
}

\begin{apient}
  bool readMemory(AddressSet::ptr addrs,
                  std::map<void *, ProcessSet::ptr> &result,
                  size_t size)
\end{apient}

\apidesc{
  This function reads memory from processes in this ProcessSet. addrs should
  contain the addresses to read memory from. size should be the amount of
  memory to read from each process. The results of the memory reads will be
  aggregated together into the result map. If any two processes read equivalent
  byte-for-byte data, then those processes are grouped together in a new
  ProcessSet associated with a common malloc allocated buffer containing their
  memory contents.
  It is the ProcControlAPI user\' s responsibility to free the
  memory buffers returned by this function.
  This function\' s behavior is undefined if addrs contains
  processes not included in this ProcessSet.
  This function returns true if every read happened without error, and false
  otherwise.
}

\begin{apient}
  bool readMemory(std::multimap<Process::const_ptr, read_t> &addr)
\end{apient}

\apidesc{
  This function reads memory from processes in this ProcessSet. The processes to
  read from are specified in the indexes of addr. The remote address, read size
  and local buffer are specified in the read\_t elements of addr.
  This function\' s behavior is undefined if addr contains
  processes not included in this ProcessSet.
  This function returns true if every read happened without error, and false
  otherwise. If any read results in an error, then the error\_ret field of the
  associated addr element will be set.
}

\begin{apient}
  bool writMemory(AddressSet::ptr addrs, const void *buffer, size_t sz) const
\end{apient}

\apidesc{
  This function will write the contents of buffer of size sz into the memory of
  each process at addrs.
  This function\' s behavior is undefined if addrs contains
  processes not included in this ProcessSet.
  This function returns true if every write happened without error, and false
  otherwise.
}

\begin{apient}
  bool writeMemory(std::multimap<Process::const_ptr, write_t> &addrs) const
\end{apient}

\apidesc{
  This function writes to the memory of each process in addrs. The processes to
  write to are specified as the indexes of addrs. The local memory buffer,
  buffer size, and target location are specified in the write\_t element of
  addrs.
  This function\' s behavior is undefined if addrs contains
  processes not included in this ProcessSet.
  This function returns true if every write happened without error, and false
  otherwise. If any write results in an error, then the error\_ret field of the
  associated addr element will be set.
}

\begin{apient}
  bool addBreakpoint(AddressSet::ptr as, Breakpoint::ptr bp) const
\end{apient}

\apidesc{
  This function inserts the Breakpoint bp into every process and
  address specified
  by as. It is similar to Process::addBreakpoint.
  This function\' s behavior is undefined if addrs contains
  processes not included in this ProcessSet.
  This function returns true if every breakpoint add happened without error, and
  false otherwise.
}

\begin{apient}
  bool rmBreakpoint(AddressSet::ptr as, Breakpoint::ptr bp) const
\end{apient}

\apidesc{
  The function removes the Breakpoint bp from each process at the locations
  specified in as. It is similar to Process::rmBreakpoint.
  This function\' s behavior is undefined if as contains
  processes not included in this ProcessSet.
  This function returns true if every breakpoint remove happened without error,
  and false otherwise.
}

\begin{apient}
  bool postIRPC(const std::multimap<Process::const_ptr, IRPC::ptr> &rpcs) const
\end{apient}

\apidesc{
  This function posts the IRPC objects specified in rpcs to their associated
  processes in the multimap. It is similar to Process::postIRPC.
  This function\' s behavior is undefined if rpcs contains
  processes not included in this ProcessSet.
  This function returns true if every post happened without error, and false
  otherwise.
}

\begin{apient}
  bool postIRPC(IRPC::ptr irpc, std::multimap<Process::ptr, IRPC::ptr> *result = NULL)
\end{apient}

\apidesc{
  This function makes a copy of irpc for each Process in this
  ProcessSet and posts
  it to that Process. If result is non-NULL, then the multimap will be filled
  with each newly created IRPC and the Process to which it was posted. It is
  similar to Process::postIRPC.
  This function returns true if every post happened without error, and false
  otherwise.
}

\begin{apient}
  bool postIRPC(IRPC::ptr irpc,
                AddressSet::ptr addrs,
                std::multimap<Process::ptr, IRPC::ptr> *result = NULL)
\end{apient}

\apidesc{
  This function makes a copy of irpc and posts it to each Process in
  addrs at the
  given Address. If result is non-NULL, then the multimap will be filled with
  each newly created IRPC and the Process to which it was posted. It is similar
  to Process::postIRPC.
  This function\' s behavior is undefined if rpcs contains
  processes not included in this ProcessSet.
  This function returns true if every post happened without error, and false
  otherwise.
}
